\documentclass[../main.tex]{subfiles}

\begin{document}

\chapter{Basics}

\section{单位分解}

\begin{theorem}{单位分解的存在性}{}
    设 \(  M  \)是一个流形, \(  \mathscr{U}= \left\{ U_{\alpha }: \alpha \in A \right\}  \)是\(  M  \)   的一个开覆盖, 则
    \begin{enumerate}
        \item 存在从属于 \(  \mathscr{U}  \)的单位分解 \(  \Phi  = \left\{  \varphi _{i}: i \in I \right\}  \)且 \(  I  \)是至多可数的, \(  \operatorname{supp} \varphi _{i}  \)是紧的. 
        \item     存在从属于 \(  \mathscr{U}  \)的具有相同指标集的单位分解 \(  \Phi = \left\{  \varphi _{\alpha } :\alpha \in A\right\}  \), 且至多可数多个 \(   \varphi _{\alpha }\neq 0  \).      
        \item 如果 \(  \mathscr{U}  \)局部有限, 且 \(  \overline{U_{\alpha }}  \)是紧的. 则存在从属于 \(  \mathscr{U}  \)的具有相同指标集的单位分解 \(  \Phi = \left\{ \phi _{\alpha }: \alpha \in A \right\}  \)    且 \(  \operatorname{supp} \varphi _{\alpha }  \)是紧的.  \footnote{由于紧急的闭子集是紧的,所以第三条可以有第二条简单地导出. }
    \end{enumerate}
    
\end{theorem}


\begin{corollary}{}{}
    设 \(  M  \)是一个流形, \(  D\subseteq M  \)是闭集,  \(  U\subseteq M  \)是 \(  D  \)的一个开邻域. 则存在 \(  f \in C^{\infty}\left( M \right)   \), 使得 \(  \left. f \right|_{D} \equiv 1  \), 且 \(  \operatorname{supp}f \subseteq U  \).       
\end{corollary}

\begin{proof}
    考虑 \(  M  \)的 一个特殊的开覆盖,  \(  \left\{ U, M\setminus D \right\}  \). 对它应用上面定理的第二种情况即可. 
\end{proof}

\section{向量场和微分算子}

\begin{definition}{}{}
\begin{enumerate}
    \item 设 \(  M  \)是一个流形, \(  p \in M  \) . 若 \(  f,g \in C^{\infty}\left( M \right)   \)满足, 存在 \(  p  \)的开邻域 \(  U  \), 使得 \(  \left. f \right|_{U} =  \left. g \right|_{U}  \), 则称 \(  f  \)和 \(  g  \)在 \(  p  \)附近等价. 
    \item 称 \(  C^{\infty}\left( M \right)/ \sim    \)为 \(  p  \)附近的函数芽, 记作 \(  C^{\infty}_{p}\left( U \right)   \) . 
    \item 对于 \(  f\in C^{\infty}_{p}\left( M \right)   \), 它的芽记作 \(  [f] \in C^{\infty}_{p}\left( M \right)   \)    
\end{enumerate}

\end{definition}

在欧式空间上, 光滑函数有展开式
\[
f\left( x \right)= f\left( x_0 \right)+ \sum _{i= 1}^{m} \frac{\partial f}{\partial x^{i}}  \left( x-x_0 \right)^{i} + \sum _{i,j= 1}^{m}\frac{1 }{2 }\frac{\partial ^{2}f}{\partial x^{i} \partial x^{j}}\left( x_0 \right)\left( x-x_0 \right)^{i}\left( x-x_0 \right)^{j}+ O\left( \left\| x-x_0 \right\|^{2} \right)      
\]

\[
    D =  C\left( x \right)+ \sum _{i= 1}^{m}a_{i}\left( x \right)\frac{\partial }{\partial x^{i}}+ \sum _{i,j= 1}^{m}\frac{1 }{2 }b_{ij}\left( x \right)\frac{\partial ^{2}}{\partial x^{i} \partial x^{j}}    
    \]

    \begin{definition}{}{}
        令\(  m_{p}= \left\{ f\in C^{\infty}_{p}\left( M \right):f\left( p \right)= 0   \right\}  \), 则 \(  m_{p}  \)是 \(  C^{\infty}_{p}\left( M \right)   \)的一个极大理想. 特别地, \(  m_{p}^{k}  \)也是 \(  C^{\infty}_{p}\left( M \right)   \)的理想.      
    \end{definition}
    则 \(  f  \)的 展开式的第二项和第三项分别视作 \(  m_{p}  \)和 \(  m_{p}^{2}  \)  中的元素.

    \begin{definition}{Differential Expression of Order \(  r  \) }{}
      \begin{enumerate}
        \item  一个 \(  r  \)阶 微分表达式, 是指一个线性函数 \(  D: C^{\infty}_{p}\left( M \right)\to \mathbb{R}    \) 满足 \(  D\left( m_{p}^{r+ 1} \right)= 0   \).  
        \item 全体 \(  p  \)处的 \(  r \)阶微分表达式记作 \(  T_{p}^{\left( r \right) }\left( M \right)   \)
        \item 定义 \(  T_{p}^{\left( \infty \right) }\left( M \right)= \bigcup _{r \ge 0} T_{p}^{\left( r \right) }\left( M \right)    \)    
      \end{enumerate}
      
    \end{definition}


    在这个定义下, 当 \(  r = 2 \)时, \[
  \begin{aligned}
    D\left( f\left( x \right)  \right)&= D\left( f\left( x_0 \right)+ \sum   \frac{\partial f}{\partial x^{i}}\left( x-x_0 \right)^{i}+ \sum \frac{1 }{2 } \frac{\partial ^{2}f}{\partial x^{i}\partial x^{j}  } \left( x-x_0 \right)^{i}\left( x-x_0 \right)^{j}   \right)   \\ 
     &= D\left( f\left( x_0 \right)  \right)+  \sum _{i} \frac{\partial f}{\partial x^{i}}\left( x_0 \right)D\left( \left( x-x_0 \right)^{i}  \right) + \sum _{i,j}\frac{1 }{2 }\frac{\partial ^{2}f    }{\partial x^{i} \partial x^{j}}   \left( x_0 \right)D\left( \left( x-x_0 \right)^{i}\left( x-x_0 \right)^{j}   \right)  \\ 
      &=: c f\left( x_0 \right)+ \sum _{i}a_{i}\frac{\partial f}{\partial x^{i}}\left( x_0 \right)+   \sum _{i,j}\frac{1 }{2 }b_{ij}\frac{\partial ^{2}f}{\partial x^{i} \partial x^{j}}\left( x_0 \right)   
  \end{aligned}
    \]其中 \(  c,a_{i},b_{i,j}  \)均为常数. 

    \begin{definition}{切向量}{}
       \begin{enumerate}
        \item  对于一个\(  1  \)阶微分表达式 \(  X  \), 若它满足 \[
        X\left( f\cdot g \right) =  X\left( f \right)\cdot g\left( x_0 \right)+ f\left( x_0 \right)X\left( g \right)     
        \]则称 \(  X  \)为 \(  p  \)处的一个切向量.    

        具体地,  \[
        X\left( f \right)= \sum _{i= 1}^{m}a_{i}\frac{\partial f}{\partial x^{i}}\left( x_0 \right)  
        \]
        \item 记 \(  T_{p}\left( M \right)   \)为 \(  p  \)处全体切向量构成的线性空间. 则 \(  T_{p}\left( M \right)\simeq \mathbb{R} ^{m}   \).  \(   \partial _{i}: =  \frac{\partial }{\partial x^{i}}  \)构成 \(  T_{p}\left( M \right)   \)的基.      
       \end{enumerate}
       
    \end{definition}
   
    让 \(  p \in U  \)在 \(  U  \)上可变 \(  \iff   \) \(  x \in V  \) 跑起来.   则 \(  X =  \sum a^{i}\left( x \right)\frac{\partial }{\partial x^{i}}   \), \(  a^{i} \in C^{\infty}\left( u \right)   \).   
    
    \begin{definition}{多重指标}{}
      \begin{itemize}
        \item   一个多重指标 \(  \alpha   \)是指 \(  \alpha = \left( \alpha _1 ,\cdots  , \alpha _{m} \right)   \) , \(  \alpha _{i} \in \mathbb{N}   \). 定义 \(  \alpha   \)的长度为 \(  \left| \alpha  \right|= \alpha _1 + \cdots + \alpha _{m}   \)
        \item 定义记号 \[
         \partial ^{\alpha } : =  \left( \frac{\partial }{\partial x^{1}} \right) ^{\alpha _1 }\cdots \left( \frac{\partial }{\partial x^{m}} \right)^{\alpha _{m}} 
        \]
      \end{itemize}
           
    \end{definition}
    
    \begin{proposition}{}{}
        \(  T_{p}^{\left( r \right) }\left( M \right)   \)的一组基为 \(  \left\{  \partial ^{\alpha }: \left| \alpha  \right|\le r  \right\}  \)  
    \end{proposition}

    \begin{definition}{}{}
        一个线性映射 $\mathcal{D}: C^\infty(M) \to C^\infty(M)$ 是一个阶数至多为 r 的微分算子，如果对于流形 $M$ 上的任意一个局部坐标图 $(U, x^1, \dots, x^m)$，算子在 $U$ 上的限制 $\mathcal{D}|_U$ 都可以表示为：
$$\mathcal{D}|_U = \sum_{|\alpha| \le r} a_\alpha(x) \partial^\alpha$$
其中：
\begin{itemize}
    \item $\alpha = (\alpha_1, \dots, \alpha_m)$ 是多重指标.
    \item $\partial^\alpha = \frac{\partial^{|\alpha|}}{(\partial x^1)^{\alpha_1} \cdots (\partial x^m)^{\alpha_m}}$ 是偏微分.
    \item 
每一个系数 $a_\alpha(x)$ 都是定义在 $U$ 上的光滑函数，即 $a_\alpha \in C^\infty(U)$.
\end{itemize}



    \end{definition}
    
    \begin{definition}{}{}
       \begin{itemize}
        \item  开集 \(  U  \)上的切向量的全体记作 \(  \mathfrak{X}\left( U \right)   \)  .  
        \item  开集 \(  U  \)上的微分算子的全体记作 \(  D\left( U\right)   \)  
       \end{itemize}
       
    \end{definition}
 \begin{remark}
    \(  D\left( U \right)   \)上按算子的复合, 可以定义出自然的乘法. 它是结合但不交换的.  
 \end{remark}

  \begin{definition}{}{}
    \(  \mathfrak{X}\left( U \right)   \)上有Lie括号 \[
   \begin{aligned}
   \left[  ,  \right]   : \mathfrak{X}\left( U \right)\times \mathfrak{X}\left( U \right)&\to \mathfrak{X}\left( U \right)\\ 
     \left[ X,Y \right]&:=  XY-YX     
   \end{aligned}
    \] 
  \end{definition}

  \begin{proposition}{}{}
    \begin{itemize}
        \item \(  C^{\infty}\left( U \right)   \)是交换结合的 \(  \mathbb{R}   \)-代数. 
        \item \(  \mathfrak{X}\left( U \right)   \)和 \(  D\left( U \right)   \)都是 \(  C^{\infty}\left( U \right)  \)模     
    \end{itemize}
    
  \end{proposition}

  \begin{proposition}{}{}
    若 \(  D_1  \)是 \(  r_1  \)阶的微分算子, \(  D_2  \)是 \(  r_2  \)阶的微分算子, 则 \begin{itemize}
        \item \(  D_1\cdot D_2  \)的阶数 \(  \le r_1+ r_2   \).        
        \item \(  D_1D_2-D_2D_1  \)的阶数 \(  \le r_1+ r_2-1  \)  
    \end{itemize}
    
  \end{proposition}

  \begin{theorem}{}{}
     设 \(  M  \)是一个流形, \(  \operatorname{dim} M = m  \) .  \(  X_1,\cdots ,X_{m}\in \mathfrak{X}\left( M \right)   \)  满足对于任意的 \(  p \in M  \). \(  X_1\left( p \right), X_2\left( p \right),\cdots , X_{m}\left( p \right)     \)构成 \(  T_{p}\left( M \right)   \)的一组基\footnote{即 \(  M  \)是一个可平行化的流形 }.    记 \(  X^{\alpha }= X_1^{\alpha _1 }\cdots X_{m}^{\alpha _{m}}  \). 则
   \begin{itemize}
    \item    \(  \forall p \in M  \) , \(  \left\{ X_{\left( p \right) }^{\alpha } :\left| \alpha  \right|\le r \right\}  \)构成 \(  T_{p}^{\left( r \right) }\left( M \right)   \)的基.   
    \item 等价地说, \(  \left\{ X^{\alpha }: \alpha = \left( \alpha _1 ,\cdots ,\alpha _{m} \right)  \right\}  \)构成 \(  D\left( M \right)   \)的   作为 \(  C^{\infty}\left( M \right)   \)模的基.  
    \item 此时, \(  D  \)可以写作 \[
    D= \sum  a_{\alpha }X^{\alpha } ,\quad a_{\alpha }\in C^{\infty}\left( M \right) 
    \] 
   \end{itemize}
   
  \end{theorem}

  \section{积分}
  
设 \(  M  \)是流形, \(  \operatorname{dim}M =  m  \), \(   \omega   \)是 \(  M  \)上的一个 \(  m  \)-形式. \(  f \in C^{\infty}\left( M \right)   \). 
  \begin{enumerate}
    \item 取 \(  M  \)的一个图册 \(  \mathcal{U}= \left\{ \left( U_{\alpha } , \varphi _{\alpha }\right): \alpha \in A  \right\}  \). 
    \item 取从属于 \(  \mathcal{U}  \)的一个紧支单位分解. \(  \Phi = \left\{  \varphi _{i}: i \in I \right\}  \)    
    \item 则 \(  f = \sum _{i \in I}f\cdot  \varphi _{i}  \). 记 \(  f_{i}= f \varphi _{i}  \)  
    \item 设 \(  \operatorname{supp} \varphi _{i}\subseteq U_{\alpha }  \) . 
    \item   设 \(  V_{\alpha }  \)中的坐标为 \(  \left(  x^1,\cdots,x^n  \right)   \) , \(  f_{i}: U_{ \alpha }\to \mathbb{R}   \). 定义 \(  \tilde{f}_{i}= f_{i}\circ  \varphi _{\alpha }^{-1} : V_{\alpha }\to \mathbb{R}   \).  
    \item  设 \(   \omega   \)在 \(  U_{\alpha }  \)上有表达式 \(   \omega  =  \omega _{\alpha }\left( x \right)\,\mathrm{d} x^{1}\wedge \cdots \wedge \,\mathrm{d} x^{m}   \)   . 
  \end{enumerate}
  \begin{definition}{对密度的积分}{}
    
  
   定义 \[
   J_{i}: =  \int_{U_{\alpha }} \tilde{f}_{i}\left( x \right)\left|  \omega _{\alpha }\left( x \right)  \right|\,\mathrm{d} x^{1}\cdots \,\mathrm{d} x^{n}  
   \]定义 \[
   J:= \int_{M}f\left|  \omega  \right|: =  \sum _{i \in I}J_{i} 
   \]称为是对密度的积分.
  \end{definition}
        
-  \begin{remark}
    需要证明 \(  J  \)与紧支单位分解的选取无关.  
  \end{remark}
  
\begin{definition}{可定向性}{}
  
  对于流形 \(  M  \). 如果存在一个图册 \(  \left\{ U_{\alpha }: \alpha \in A \right\}  \)使得任意两个坐标卡的之间的转移函数的Jacobi行列式都是正的. 则称 \(  M  \)是\dfntxt{可定向}的.   
  \begin{itemize}
    \item  称这样的图册为\dfntxt{定向图册}
    \item 若 \(  \mathcal{U_1}  \)和 \(  \mathcal{U_2}  \)是两个定向图册, 满足 \(  \mathcal{U}_{1}\cup  \mathcal{U}_{2}  \)还是定向图册,则称 \(  \mathcal{U}_{1}  \)和 \(  \mathcal{U}_{2}  \)是\dfntxt{定向相容}的.     
  \end{itemize}
  
\end{definition}

\begin{definition}{对形式的积分}{}
若\(  M  \)还是可定向的,可以  定义 \[
  J_{i}= \int_{U_{\alpha }} \tilde{f}_{i}\left( x \right) \omega _{\alpha }\left( x \right)\,\mathrm{d} x^{1}\cdots \,\mathrm{d} x^{n}  
  \]定义 \[
  J: = \int_{M}f  \omega : =   \sum _{\in I}J_{i}
  \] 
\end{definition}
\begin{remark}
  可以证明, \(  J  \)不依赖于与 \(  \mathcal{U}  \)定向相同的图册的选取.   
\end{remark}

\begin{definition}{}{}
  \(  M  \)是可定向的, 当且仅当 \(  M  \)上存 在一个处处非零的 \(  m  \)- 形式.    
\end{definition}

\begin{theorem}{}{}
  设 \(  M  \)是可定向的带边流形, \(   \omega   \)是一个紧支的 \(  \left( m-1 \right)   \)   形式, 则 \[
  \int_{M}\,\mathrm{d}  \omega = \int_{ \partial M} \omega 
  \]
\end{theorem}

\section{映射与子流形}

\begin{lemma}{}{}
    设 \(  M,N  \)是 \(  C^{\infty}  \)流形, \(   \varphi :M\to N \)连续   . 若对于任意的\(  N  \)的开集 \(  V  \), 和 \(  g \in C^{\infty}\left( V \right)   \). \(  g\circ  \varphi \in C^{\infty}\left(  \varphi^{-1}\left( V \right)  \right)  \) . 则 \(   \varphi  \)是 \(  C^{\infty}  \)的.       
\end{lemma}
 
\begin{lemma}{}{}
    对于 \(  y \in N  \)附近的两个 \(  C^{\infty}  \)函数 \(  g_1,g_2  \), 若 \(  g_1\sim_{y}g_2  \)(函数芽) . 则 \(  g_1\circ  \varphi\sim_{x}g_2\circ  \varphi  \).     
\end{lemma}  

\begin{definition}{拉回映射}{}
    \(  M,N  , \varphi\)同前 . 则 \(   \varphi  \)诱导出线性  映射 \[
     \begin{aligned}
     \varphi ^{*}: C^{\infty}_{y}\left( N \right)&\to C^{\infty}_{x}\left( M \right)  \\ 
      g\mapsto  g\circ  \varphi
     \end{aligned}
    \]称为是关于 \(   \varphi  \)的\dfntxt{拉回映射} 
\end{definition}

\begin{definition}{切映射}{}
     \(  M,N  , \varphi\)同前. 则 \(   \varphi  \)诱导出对偶空间(微分算子空间)之间的映射 \[
     \left( \,\mathrm{d} \varphi \right) _{*} : C^{\infty}_{x}\left( M \right)^{*}\to C^{\infty}_{ \varphi\left( x \right) }\left( N \right)^{*}  
     \]  通过做限制得到映射 
    \begin{enumerate}
      \item \[
      \left( \,\mathrm{d} \varphi \right)_{x}: T_{x}M\to T_{ \varphi\left( x \right) }N 
      \]
      \item \[
      \left( \,\mathrm{d} \varphi \right)_{x}^{\left( \infty \right) } : C^{\infty}_{x}\left( M \right)^{*}\to C^{\infty}_{ \varphi\left( x \right) }\left( N \right)^{*}  
      \]
    \end{enumerate}
    
\end{definition}
\begin{remark}
    这里蕴含了一些事实
    \begin{enumerate}
      \item 如果 \(  X  \)是切向量, 则 \(  \left( \,\mathrm{d} \varphi \right)_{x}\left( X \right)    \)也是.   
      \item \(   \varphi^{-1*}\left( m_{y} \right)\subseteq m_{x}   \) , 从而 \(   \varphi^{-1*}\left( m_{y}^{k} \right)\subseteq m_{x}^{k}   \). 
      \item 若 \(  D  \)是一个 \(  r  \)阶微分表达式, \(  \left( \,\mathrm{d} \varphi \right)_{x}\left( D \right)    \)     也是 \(  r  \)阶微分表达式.  
    \end{enumerate}
    
\end{remark}

\begin{definition}{}{}
    
\(  M,N, \varphi  \)同前. 若 \(  X  \)是 \(  M  \)上的向量场.
    若 \(  X \in \mathfrak{X}\left( M \right)   \), \(  Y \in \mathfrak{X}\left( N \right)   \)满足 \[
    \forall_{x \in M} \left( \,\mathrm{d} \varphi \right)_{x}\left( X\left( x \right)  \right)=Y\left(  \varphi\left( x \right)  \right)   
    \]  则称 \(  X  \)和 \(  Y  \)是 \(   \varphi  \)   -相关的. 
\end{definition}

\begin{lemma}{}{}
     设 \(  D  \)是 \(  M  \)上的微分算子. \(  D: C^{\infty}\left( M \right)\to  C^{\infty}\left( M \right)    \) . 若 \(  \tilde{D}\in D\left( C^{\infty}\left( N \right)  \right)  \) . 则 \(  \tilde{D}: C^{\infty}\left( N \right)\to C^{\infty}\left( N \right)    \)
\end{lemma}

\begin{definition}{}{}
    \(   \varphi:M\to N  \)诱导出映射 \[
     \varphi ^{*}:  \Omega^{*}\left( N \right)\to \Omega^{*}\left( M \right)  
    \] 
\end{definition}


\begin{definition}{}{}
    设 \(  M,N , \varphi  \)同前. \(  \operatorname{dim}M=m, \operatorname{dim}=n  \). \(  r_{x}: =\operatorname{rank}\left( \,\mathrm{d} \varphi \right)_{x}   \). 
    \begin{enumerate}
      \item 若对于任意的 \(  x \in M  \), \(  r_{x}=n  \). 则称 \(   \varphi  \)是一个淹没(Submersion). 
      \item 若对于任意的 \(  x \in M  \) , \(  r_{x} =m  \) . 则称 \(   \varphi  \)是一个浸入(immersion) .       
    \end{enumerate}
       
\end{definition}

\begin{proposition}{}{}
    \begin{enumerate}
      \item 若 \(   \varphi  \)是一个淹没, 则 \(  m\ge n  \). 且存在局部坐标系, 使得 \[
     \tilde{\varphi}\left(  x^1,\cdots,x^m \right)=\left(  x^1,\cdots,x^n \right)  
    \]
    \item 若 \(   \varphi  \)是一个浸入, 则 \(  m\le n  \). 且存在局部坐标, 使得 \[
     \tilde{\varphi}\left(  x^1,\cdots,x^m \right)=\left(  x^1,\cdots,x^n,0,\cdots,0 \right)  
    \]  
    \end{enumerate}
      
\end{proposition}


\begin{definition}{}{}
    设 \(  M,N  \)是拓扑流形. \(   \varphi:M\to N  \)是连续映射. 
    \begin{enumerate}
      \item 若对于任意的 \(  x \in M  \). 记 \(  y= \varphi\left( x \right)   \) . 存在 \(  x,y  \)附近的坐标卡 \(  \left( U, \varphi_{U}: U\to \tilde{U} \right), (V ,  \varphi_{v}:V\to \tilde{V})   \).使得\[
       \tilde{\varphi}= \varphi_{V}\circ  \varphi  \circ  \varphi_{U}^{-1}: \tilde{U}\to \tilde{V}
      \]是一个从 \(  \mathbb{R}^{m}  \)到 \(  \mathbb{R}^{n}  \)的投影. 则称 \(   \varphi  \)是一个拓扑淹没.         
      \item 可以类似地定义出拓扑浸入. 
    \end{enumerate}
      
\end{definition}

\begin{theorem}{}{}
   设 \(  M  \)设 \(  C^{\infty}  \)流形. \(  N  \)是拓扑流形. 若 \(   \varphi: M\to N  \)是拓扑的淹没且满, 则存在唯一的 \(  C^{\infty}  \)结构, 使得 \(   \varphi  \)是一个 \(  C^{\infty}  \)淹没.  
\end{theorem}

\begin{definition}{}{}
  设 \(   \varphi   :M\to N\)是一个浸入. 
  \begin{enumerate}
    \item 若 \(   \varphi   \)是单的, 则称 \(   \varphi   \)是一个\dfntxt{嵌入}. 
    \item 若 \(   \varphi   \)是嵌入, 且 \(   \varphi _{M}: M\to  \varphi \left( M \right)   \)    是同胚. 则称 \(   \varphi   \)是一个\dfntxt{正则嵌入. } 
  \end{enumerate}
   
\end{definition}

\begin{theorem}{正则嵌入的水平集刻画}{}
  设 \(   \varphi :M\to N  \)是正则嵌入. 则对于 \(  y \in  \varphi \left( M \right)   \). 存在 \(  y  \)在 \(  N  \)中的开邻域 \(  V  \), 以及 \(  V  \)上的\(  C^{\infty}  \)映射 \(  F: V\to \mathbb{R}   ^{n-m}\). 使得  \(   \varphi \left( M \right)\cap V= \left\{ z\in V: F\left( z \right)= 0  \right\}   \)          
\end{theorem}

\begin{theorem}{正则嵌入的参数刻画}{}
  对于 \(  x \in M , y =  \varphi \left( x \right)\in  \varphi \left( M \right)   \). 存在 \(  x  \)在 \(  M  \)中的开邻域 \(  U  \), \(  y  \)在 \(  N  \)中的开邻域 \(  V  \). 使得 \(   \varphi \left( U \right)=  \varphi \left( M \right)\cap V    \).         
\end{theorem}

\begin{theorem}{}{}
  设 \(  P  \)是任一 \(  C^{\infty}  \)流形.  \(   \varphi :M\to N  \)是正则嵌入.   \(  u: P\to M  \)是任一映射. 则 \[
 u \text{是} C^{\infty}\text{的} \iff   \varphi \circ u:P\to N \text{是} C^{\infty} \text{的}. 
  \]   
\end{theorem}

\begin{definition}{}{}
  若 \(   \varphi :M\to N  \)是嵌入且满足\[
 u \text{是} C^{\infty}\text{的} \iff   \varphi \circ u:P\to N \text{是} C^{\infty} \text{的}. 
  \]    则称 \(   \varphi   \)是\dfntxt{拟正则嵌入} .   
\end{definition}

\begin{definition}{}{}
  设 \(  N  \)是一个 \(  C^{\infty}  \)流形.  \(  M\subseteq N  \). 若 \(  i: M\to N  \)是嵌入. 则称 \(  M  \)是 \(  N  \)的子流形. 若 \(  i  \)还是正则或拟正则的. 则称 \(  M  \)是 \(  N  \)的正则或拟正则子流形.          
\end{definition}

\begin{theorem}{正则值}{}
  设 \(  M,N  \)是 \(  C^{\infty}  \)流形. \(   \varphi :M\to N  \)是 \(  C^{\infty}  \)映射. 对于 \(  y \in N  \). 记 \(  M^{\prime} =  \varphi ^{-1} \left( y \right)   \) . 若对任一的 \(  x \in M^{\prime}   \). \(  \operatorname{rank}\left( \,\mathrm{d}  \varphi  \right)_{x}= n   \). 则 \(  M^{\prime}   \)是 \(  M  \)的一个正则子流形, 维数为 \(  m-n  \).   这样的 \(  y  \)叫 \(   \varphi   \)的一个正则值.            
\end{theorem}
\begin{note}
  \begin{itemize}
    \item \(  m-n  \)衡量的是 \(  M  \)上 有多少维数的信息, 在通过\(   \varphi   \)映射的过程中被压缩掉了. 
    \item  \(  y  \)的正则性, 体现在它的“来源信息”在 \(  M  \)上是光滑的. 
    \item 正则值 \(  y  \)给出的正则子流形 \(  M^{\prime}   \), 就是 \(  M  \)上那些在映射 \(   \varphi   \)下被压缩成 \(  y  \)的信息构成的子流形.        
  \end{itemize}
  
\end{note}
\begin{example}{}{}
  \begin{enumerate}
    \item \(  \mathrm{Mat}\left( n,\mathbb{R}  \right)\simeq \mathbb{R} ^{2}   \)是\(  n^{2}  \)维流形.  
    \item \(  \operatorname{GL} \left( n,\mathbb{R}  \right)= \left\{ M\in \mathrm{Mat}\left( n,\mathbb{R}  \right):\det M\neq 0  \right\}   \)作为 \(  \mathrm{Mat}\left( n,\mathbb{R}  \right)   \)的一个开子集也是一个\(  n^{2}  \)维流形.    
    \item  \(  \mathrm{SL}\left( n,\mathbb{R}  \right)= \left( M\in \mathrm{Mat}\left( n,\mathbb{R}  \right)  \right):\det M= 1    \)   \[
    \begin{aligned}
    \det : \mathrm{Mat}\left( n,\mathbb{R}  \right)&\to \mathbb{R} \\ 
     M&\mapsto \det M  
    \end{aligned}
    \]  \(  SL  \left( n,\mathbb{R}  \right)= \det ^{-1} \left( I \right)  \). 取一个 \(  M \in \mathrm{Mat}\left( n,\mathbb{R}  \right)   \), \(  \det M= 1  \). 要证明 \(  \left( \,\mathrm{d} \left( \det  \right)  \right)_{M}: \mathbb{R} ^{n^{2}}\to \mathbb{R}    \)   . 按分量求偏导,得到 \[
    \frac{\partial }{\partial a_{ij}} \left( \det M \right)=  A_{ij}(\text{代数余子式})  
    \]由于 \(  \det M= 1  \).  至少有一个 \(  A_{ij}\neq 0  \). 从而  \(  \left( \,\mathrm{d} \left( \det  \right)  \right)_{M}   \)是满的. 
  \item 定义 \(  P: \mathrm{Mat}\left( n,\mathbb{R}  \right)\to  \mathrm{SMat}\left( n,\mathbb{R}  \right)(\text{对称矩阵}) \simeq \mathbb{R} ^{\frac{n\left( n+ 1 \right)  }{2 } }   \) .  \[
  P\left( A \right): = A^{\top}A 
  \]则 \(  O\left( n,\mathbb{R}  \right)= P^{-1} \left( I \right)   \) . 对于 \(  A \in O\left( n,\mathbb{R}  \right)   \). \[
  \left( \,\mathrm{d} P \right)_{A}: \mathrm{Mat}\left( n,\mathbb{R}  \right)\to \mathrm{SMat}\left( n,\mathbb{R}  \right)   
  \]与上例不同, 这里用分量的形式去计算 Jacobi是看不清楚的, 不过由于 \(  P  \)本身足够简单, 我们可以用微分最原始的含义来观察. 考虑  \[
  P\left( A+  \varepsilon B \right)-P\left( A \right)=  \varepsilon \left( \,\mathrm{d} P \right)_{A}\left( B \right)+ O\left(  \cdots   \right)     
  \]带入得到 \[
  \left( A^{T}+  \varepsilon B^{T} \right)\left( A+  \varepsilon B \right)-A^{\top}A=    \varepsilon \left(B^{\top}A+ A^{\top}B \right)+ O\left(  \varepsilon  \right)  
  \]于是 \(  \left( \,\mathrm{d} P \right)_{A}   \left( B \right)= B ^{\top}A+ A^{\top}B\) . 只需要证明这个关于 \(  B  \)的映射是满的. 为此, 解方程 \[
  A^{\top}B+ B^{\top}A= C
  \] 取 \(  B = \frac{1}{2}AC  \), 利用 \(  A,C  \)对称,   带入得到 \[
  \frac{1}{2}A^{\top}AC+ \frac{1}{2}C^{\top}A^{\top}A= \frac{1}{2}\left( C+ C^{\top} \right)= C 
  \]  
  故 \(  O\left( n,\mathbb{R}  \right)   \)是 \(  \mathrm{Mat}\left( n,\mathbb{R}  \right)   \)的一个 \(  \frac{n\left( n-1 \right)  }{2 }   \)维正则子流形. 注意到对于 \(  A \in O\left( n,\mathbb{R}  \right)   \) , \(  \det A= 1  \)或 \(  \det A= -1  \) , 可以窥见事实上 \(  O\left( n,\mathbb{R}  \right)   \)有两个连通分支.  
  \end{enumerate}
  
\end{example}

\section{Frobenius定理}
\begin{definition}{\(  C^{\infty}  \)分布 }{}
  在光滑流形 $M$ 上, 一个 $k$-维 $C^{\infty}$ 分布（smooth distribution）是一个映射 $\Delta: M \to \bigcup_{p \in M} \operatorname{Gr}(k, T_p M)$, 这里\[
  \operatorname{Gr}(k, V) = \{ W \subseteq V \mid W \text{ 是 } k \text{-维子空间} \}
  \] 其中: 
\begin{itemize}
  \item 对于每个 $p \in M$, $\Delta_p$ 是切空间 $T_p M$ 的一个 $k$-维子空间.
  \item $\Delta$ 是光滑的,即它对应于切丛 $TM$ 的一个光滑子向量丛: 局部存在光滑向量场 $X_1, \dots, X_k$ 生成 $\Delta_p = \operatorname{span}\{X_1(p), \dots, X_k(p)\}$.
\end{itemize}

\end{definition}

\begin{definition}{}{}
     设 \(  \mathscr{L}  \)是 \(  N \)上的一个分布.   
  \begin{itemize}
    \item 设 \(\mathscr{L} \) 是一个 \(  m  \)维分布.  \(  M  \)是 \(  N  \)的一个 \(  m  \)维子流形. 若对于 \(  x \in M  \) . \(  T_{x}M= L_{x}  \), 则称 \(  M  \)是 \(  \mathscr{L} \)的一个积分子流形.          
    \item 若对于任意的 \(  x \in N  \) . 都存在一个过 \(  x  \)的 \(  \mathscr{L}  \)的积分子流形. 则称 \(  \mathscr{L}  \)是可积的.    
    \item 若 \(  \left[ \mathscr{L},\mathscr{L} \right]\subseteq \mathscr{L}   \), 则称 \(  \mathscr{L}  \)是对合的 (involutive) .    
  \end{itemize}
  
\end{definition}
\begin{remark}
  由于Lie bracket是封闭的, 从而可积一定是对合的.
\end{remark}


\begin{theorem}{局部Frobenius}{}
  设 \(  \mathscr{L}  \)是 \(  N  \)上的一个分布.  若 \(  \mathscr{L}  \)是对合的, 则 \(  \mathscr{L}  \)是可积的.     
\end{theorem}

\begin{theorem}{整体Frobenius}{}
  设 \(  N  \)是一个 \(  n  \)维  \(  C^{\infty}  \)流形. \(  \mathscr{L}  \)是 \(  N  \)上的一个 \(  m  \)维对合分布  . 则过 \(  N  \)上任一点 \(  p  \) , 存在唯一的 \(  \mathscr{L}  \)的极大积分子流形 \(  M_{p}  \) .        \(  \mathscr{L}  \)的任何积分子流形都是 \(  N  \)的拟正则子流形, 并且是某一个极大积分子流形的开子流形. 
\end{theorem}
\begin{proofsketch}
  对 \(  p \in N  \), \[
  M_{p}= \left\{ q \in N: \text{存在一条分片光滑的 }\mathscr{L}\text{的积分曲线, 连接}p\text{和}q \right\}
  \] 参考Warner的GTM94
\end{proofsketch}
\end{document}